\section*{Chapter \ref{ch:b2:ore-deposits} --- Ore Deposits}

\begin{solution}{exo:b2:ore-deposits:1}
(a) $\qty{0.5}{\percent} = \qty{5000}{ppm}$, and $5000/28 = 179$. (b) $\qty{6}{g/t} = \qty{6000}{ng/g}$,
and $6000/1.5 = 4000$. (c) Chromite holds $2 \times 52.00/(55.85 + 104.00 + 64.00) = 0.4646$ of
chromium by mass; the chromitite, $0.9 \times 0.4646 = 0.418$, or \qty{418000}{ppm}, and
$418000/92 = 4540$.
\end{solution}

\begin{solution}{exo:b2:ore-deposits:2}
(a) $914 \times 0.0053 = \qty{4.84}{Mt}$ of copper and $914 \times 10^6 \times 0.40 = \num{3.66e8}$ g,
\qty{366}{t} of gold; $1640 \times 0.0044 = \qty{7.22}{Mt}$ and \qty{738}{t}. (b) $1640 - 914 =
\qty{726}{Mt}$ holding $7.216 - 4.844 = \qty{2.37}{Mt}$ of copper, a mean grade of
$2.372/726 = \qty{0.327}{\percent}$: the extra blocks lie below the \qty{0.3}{\percent} cut-off
and just above it. (c) They answer the question at two cut-offs; neither is the size of the
deposit without its cut-off.
\end{solution}

\begin{solution}{exo:b2:ore-deposits:3}
(a) $(19.3 - 1.00)/(2.65 - 1.00) = 18.3/1.65 = 11.1$. (b) Equal speeds need $(\rho_g - 1)d_g^2 =
(\rho_q - 1)d_q^2$, so $d_g/d_q = \sqrt{1.65/18.3} = 0.300$: a gold grain settles like a quartz
grain about three times larger. (c) $d_q = 0.1\sqrt{3.7/1.65} = \qty{0.150}{mm}$.
\end{solution}

\begin{solution}{exo:b2:ore-deposits:4}
Per unit area and density, the blanket keeps its own copper and receives that of the leached
column: $g = g_0(h_e + \eta h_l)/h_e$. (a) $0.4 \times (50 + 150)/50 = \qty{1.6}{\percent}$.
(b) $0.4 \times (50 + 0.8 \times 150)/50 = \qty{1.36}{\percent}$. (c) $h_l = 50 \times (1.7/0.4 - 1)
= \qty{163}{m}$. These grades sit in the 0.4--\qty{1.7}{\percent} range of enriched porphyry
blankets.
\end{solution}

\begin{solution}{exo:b2:ore-deposits:5}
(a) Platinum, $C_S = 10 \times 10^4 (R + 1)/(R + 10^4)$ ppb: \qty{1.00}{ppm} for $R = 100$,
\qty{9.10}{ppm} for \num{1000}, \qty{50.0}{ppm} for \num{10000}, \qty{90.9}{ppm} for \num{100000}.
Nickel, $C_S = 200 \times 500 (R + 1)/(R + 500)$ ppm: \qty{1.68}{\percent}, \qty{6.67}{\percent},
\qty{9.52}{\percent}, \qty{9.95}{\percent}. (b) $D\,C_L$: \qty{100}{ppm} of platinum,
\qty{10}{\percent} of nickel. (c) Nickel is within \qty{5}{\percent} of its ceiling from
$R = \num{10000}$ on, so its content hardly changes over a factor of ten in $R$; platinum is still
nearly doubling there, so its content measures $R$.
\end{solution}

\begin{solution}{exo:b2:ore-deposits:6}
(a) $\mu_M = 2.373 - 0.364 - 2 = 0.009$; $\sigma_M = 0.702$; median $10^{0.009} = \qty{1.02}{Mt}$.
(b) $z = (1 - 0.009)/0.702 = 1.412$, $1 - \Phi(1.412) = 0.0790$: about 33 deposits of 422.
(c) $\sigma = \sqrt{0.654^2 + 0.199^2} = 0.684$, $z = 1.449$, $1 - \Phi = 0.0737$, about 31: the
weak correlation adds a few giants.
\end{solution}

\begin{solution}{exo:b2:ore-deposits:7}
Mean of all blocks $0.40\,e^{0.18} = \qty{0.479}{\percent}$. For $c = 0.3$: $z = \ln(0.40/0.3)/0.60 =
0.4795$, $\Phi(z) = 0.6842$, $\Phi(z + s) = 0.8598$: $T = \qty{684}{Mt}$, $\bar g = 0.479 \times
0.8598/0.6842 = \qty{0.602}{\percent}$, $M = 1000 \times 0.00479 \times 0.8598 = \qty{4.12}{Mt}$.
For $c = 0.6$: $z = -0.6758$, $\Phi = 0.2496$, $\Phi(z + s) = 0.4698$: \qty{250}{Mt} at
\qty{0.901}{\percent}, \qty{2.25}{Mt}. On the figure, at a cut-off of \qty{0.3}{\percent} the blue
curve reads about \qty{680}{Mt} and the red one about \qty{0.6}{\percent}.
\end{solution}

\begin{solution}{exo:b2:ore-deposits:8}
(a) Fluid: $\num{5e6}/(\num{1e-3} \times 0.5) = \qty{1.00e10}{t}$; \omterm{def:b2:magma-emplacement:magma}{magma}: $\num{1e10}/0.035 =
\qty{2.86e11}{t}$. (b) $\num{2.86e14}/2600 = \qty{1.10e11}{m^3}$, \qty{110}{km^3}, a cube of \omterm{def:b2:magma-emplacement:magma}{magma}
about \qty{4.8}{km} on a side. (c) The \omterm{def:b2:magma-emplacement:magma}{magma} held $\num{2.86e11}
\times \num{50e-6} = \qty{1.43e7}{t}$ of copper, and the ore $5/14.3 = 0.350$ of it.
\end{solution}

\begin{solution}{exo:b2:ore-deposits:9}
(a) Mean of the logarithms 1.411, standard deviation 0.0664; median $10^{1.411} = \qty{25.8}{ppm}$;
threshold $10^{1.411 + 0.133} = 10^{1.544} = \qty{35.0}{ppm}$. (b) $(4 \times 400 + 36 \times 25.8)/40 =
\qty{63.2}{ppm}$, above \qty{35.0}{ppm}: still anomalous after dilution by nine parts of clean
sediment. (c) $4 \times 400 + (A - 4) \times 25.8 = 35.0\,A$ gives $A = 4 \times (400 -
25.8)/(35.0 - 25.8) = \qty{163}{km^2}$.
\end{solution}

\begin{solution}{exo:b2:ore-deposits:10}
(a) $e^{su}\,e^{-u^2/2} = e^{-(u - s)^2/2 + s^2/2}$; divide by $\sqrt{2\pi}$. (b) $E[e^{\mu + su}] =
e^{\mu}\int e^{su}\varphi(u)\,\mathrm{d}u = e^{\mu + s^2/2}\int\varphi(u - s)\,\mathrm{d}u = e^{\mu +
s^2/2}$. (c) $\Phi$ increases and $s > 0$, so $\Phi(z + s) > \Phi(z)$: the blocks kept are richer than
the blocks discarded, so each cut-off keeps a larger share of the metal than of the rock; the
difference grows with $s$, the variability of the deposit.
\end{solution}

\begin{solution}{exo:b2:ore-deposits:11}
(a) $D(R + 1)/(R + D) = fD$ gives $R + 1 = fR + fD$, $R = (fD - 1)/(1 - f)$. (b) $f = 0.5$: $R = D - 2$,
298 and \num{9998}; $f = 0.9$: $R = 9D - 10$, \num{2690} and \num{89990}. (c) Nickel reaches
\qty{50}{\percent} of its ceiling inside the range of nickel--copper ores (100 to \num{2000}), and 90 just above it;
platinum needs $R$ of \num{10000} to \num{100000}, the range of the PGE reefs.
\end{solution}

\begin{solution}{exo:b2:ore-deposits:12}
(a) With $\ln M = m + s u$, the deposits above the quantile have $u > \Phi^{-1}(q)$; their metal is
$\int_{\Phi^{-1}(q)}^{\infty} e^{m + su}\varphi(u)\,\mathrm{d}u = e^{m + s^2/2}[1 - \Phi(\Phi^{-1}(q) -
s)]$, and dividing by the total $e^{m + s^2/2}$ gives $\Phi(s - \Phi^{-1}(q))$. (b) $s = 0.702 \ln 10 =
1.616$; $\Phi(1.616 - 1.282) = \Phi(0.334) = 0.631$. (c) $\Phi(1.616 - 2.326) = \Phi(-0.710) = 0.239$,
against 0.266 for the four largest: the largest deposits are a little larger than the log-normal
law expects.
\end{solution}

\begin{solution}{pb:b2:ore-deposits:1}
\textbf{1.} \qty{330}{m}; $\sum l\,g = 138.4$, mean $138.4/330 = \qty{0.419}{\percent}$.
\textbf{2.} $3.24/8 = \qty{0.405}{\percent}$; each interval stands for a length of core, hence a
volume of rock: a long interval must count more than a short one.
\textbf{3.} $230/330 = 0.697$, at $118.6/230 = \qty{0.516}{\percent}$.
\textbf{4.} $4190/28 = 150$.
\textbf{5.} $8100/28 = 289$.
\textbf{6.} $0.38\,e^{0.55^2/2} = 0.38\,e^{0.1513} = \qty{0.442}{\percent}$; the hole, at
\qty{0.419}{\percent}, is within \qty{6}{\percent} of it: representative.
\textbf{7.} $\pi \times 450^2 \times 700 = \qty{4.45e8}{m^3}$; $\times 2.7 = \num{1.20e9}$ t,
\qty{1200}{Mt}.
\textbf{8.} $z = \ln(0.38/0.3)/0.55 = 0.4298$; $\Phi(z) = 0.6663$.
\textbf{9.} $z = 1.167$, 0.4298, $-0.4990$; $\Phi = 0.8784$, 0.6663, 0.3089; $T = \qty{1050}{Mt}$,
\qty{800}{Mt}, \qty{371}{Mt}.
\textbf{10.} $\Phi(z + s) = \Phi(0.9798) = 0.8364$; $\bar g = 0.442 \times 0.8364/0.6663 =
\qty{0.555}{\percent}$.
\textbf{11.} $0.442 \times 0.9570/0.8784 = \qty{0.482}{\percent}$; $0.442 \times 0.5203/0.3089 =
\qty{0.744}{\percent}$.
\textbf{12.} \omterm{def:b2:ore-deposits:grade}{Tonnage} $\times\, 370.7/799.6 = 0.464$; grade $\times\, 0.744/0.555 = 1.34$.
\textbf{13.} At the median, \qty{0.38}{\percent} ($z = 0$): \qty{600}{Mt} at $0.442 \times
\Phi(0.55)/0.5 = 0.442 \times 0.7088/0.5 = \qty{0.627}{\percent}$.
\textbf{14.} $1200 \times 0.00442 = \qty{5.30}{Mt}$.
\textbf{15.} $(\log_{10} 1200 - 2.373)/0.654 = (3.0792 - 2.373)/0.654 = 1.080$; $\Phi = 0.8599$: the
86th percentile, a large deposit.
\textbf{16.} $(\log_{10} 0.442 + 0.364)/0.199 = (-0.3546 + 0.364)/0.199 = 0.047$; $\Phi = 0.5187$: an
ordinary grade, 52nd percentile.
\textbf{17.} 0.702 with the correlation, 0.684 without.
\textbf{18.} $\mu_M = 0.009$; $(\log_{10} 5.30 - 0.009)/0.702 = (0.7243 - 0.009)/0.702 = 1.019$; $\Phi =
0.8459$: the 85th percentile.
\textbf{19.} $422 \times (1 - 0.8459) = 65$.
\textbf{20.} $M = T_0\,e^{\mu + s^2/2}\,\Phi(z + s) = 1200 \times 0.00442 \times 0.9570 =
\qty{5.08}{Mt}$.
\textbf{21.} $1200 \times 0.00442 \times 0.5203 = \qty{2.76}{Mt}$.
\textbf{22.} $\Phi(z + s) = 0.8364$ of the copper, in $\Phi(z) = 0.6663$ of the \omterm{def:b2:ore-deposits:grade}{tonnage}.
\textbf{23.} From 0.3 to \qty{0.2}{\percent} the \omterm{def:b2:ore-deposits:grade}{tonnage} grows by $1054 - 800 = \qty{254}{Mt}$
(\qty{32}{\percent}) and the copper by $5.08 - 4.44 = \qty{0.64}{Mt}$ (\qty{14}{\percent}): the added
blocks average $0.64/254 = \qty{0.25}{\percent}$, the poorest of the deposit.
\textbf{24.} $1200 \times 0.00442 \times 0.8364 = 800 \times 0.00555$: the prospect contains
\textbf{\qty[text-series-to-math=true]{4.44}{Mt}} of copper above the \qty{0.3}{\percent} cut-off.
\end{solution}
